% !PRINTED 500
\weberpage{500}
%<<CONT
wenn $\varphi\psi_1 = \varphi_1\psi$ ist. Haben die beiden Functionen $\varphi, \psi$ keinen gemeinsamen Theiler, so heisst $\varphi:\psi$ ein \gesperrt{irreducibler Bruch}. Unter den verschiedenen Darstellungen einer Function $\omega$ giebt es \gesperrt{eine} durch einen irreduciblen Bruch, und diese wollen wir die \gesperrt{einfachste Darstellung} nennen. In ihr sind Z\"ahler und Nenner bis auf einen gemeinsamen Factor, der eine beliebige Zahl in $\Omega$ sein kann, durch $\omega$ selbst v\"ollig bestimmt (Bd. I, \S.\,141).

Eine solche Function $\omega$ wollen wir ein \gesperrt{Functional des K\"orpers} $\Omega$ nennen. Als specielle F\"alle sind darunter die ganzen Functionen und die Zahlen selbst enthalten. Bei den Zahlen ist \gesperrt{jede} Darstellung als Quotient zweier Zahlen in $\Omega$ als einfachste Darstellung zu betrachten.

Auf die Functionale lassen sich die vier Grundrechnungsarten in demselben Umfange anwenden, wie auf die Zahlen, und der Inbegriff aller Functionale des K\"orpers $\Omega$ ist daher gleichfalls ein K\"orper, den wir mit $\overline{\Omega}$ bezeichnen und den \gesperrt{Functionalk\"orper} von $\Omega$ nennen wollen.\footnote{Es liegt nahe, die Functionale des K\"orpers $\Omega$, mit denen wie mit den Zahlen in $\Omega$ gerechnet wird, geradezu als \gesperrt{ideale Zahlen} des K\"orpers $\Omega$ zu bezeichnen. Diese Ausdrucksweise w\"urde sich einerseits an die Idealfactoren von \gesperrt{Kummer}, andererseits an die von \gesperrt{Dedekind} eingef\"uhrten Ideale anschliessen, die zu den Functionalen in einer sehr nahen, sp\"ater zu er\"orternden Beziehung stehen. So bestechend in mancher Hinsicht eine solche Terminologie w\"are, so erweist sie sich doch in anderer Hinsicht nicht als zweckm\"assig, weil dadurch den Functionalen, die doch immerhin nur Mittel zum Zweck, nicht selbst\"andig f\"ur sich Gegenstand unseres Interesses sind, anscheinend eine gr\"ossere Bedeutung beigelegt w\"urde, als ihnen in Wirklichkeit zukommt. Mit der Einf\"uhrung des Wortes \glqq Functional\grqq{} folgen wir einem Vorschlage von \gesperrt{Dedekind}.}

Der Functionalk\"orper $\overline{\Omega}$ enth\"alt den Zahlk\"orper $\Omega$ als Theiler.

Ein Functional, dessen Co\"efficienten rationale Zahlen sind, heisst ein \gesperrt{rationales Functional}. Der Inbegriff aller rationalen Functionale ist der Functionalk\"orper $\overline{R}$ des K\"orpers $R$ der rationalen Zahlen, und der K\"orper $\overline{R}$ ist in jedem algebraischen Functionalk\"orper $\overline{\Omega}$ enthalten.

Jedes rationale Functional $A$ kann als Quotient zweier ganzer Functionen in $R$ dargestellt werden, denen man auch ganzzahlige Co\"efficienten geben kann, indem man Z\"ahler und Nenner mit dem Hauptnenner aller Co\"efficienten multiplicirt.
